\subsection{Subspaces and Spans}

\begin{definition} \label{def:subspace} \hlt{Subspace}\\
\lean{MathNotes.LinearAlgebra.StructuralLinearAlgebra.VectorSpacesAndLinearStructure.Subspace}\leanok Let $\mathcal{V}$ be a vector space over $\mathbb{F}$.\\
A \hlt{subspace} of $\mathcal{V}$ is a subset $\mathcal{U}\subseteq\mathcal{V}$ that is an $\mathbb{F}$-vector space under the inherited operations; write $\mathcal{U}\leq\mathcal{V}$.
\end{definition}

\begin{theorem} \label{thm:subspacetest} \hlt{Subspace Test}\\
\lean{MathNotes.LinearAlgebra.StructuralLinearAlgebra.VectorSpacesAndLinearStructure.subspace_test}\leanok Let $\mathcal{V}$ be a vector space over $\mathbb{F}$ and let $\emptyset\neq\mathcal{U}\subseteq\mathcal{V}$. Then $\mathcal{U}\leq\mathcal{V}\iff(\forall\,\mathbf{u},\mathbf{v}\in\mathcal{U})(\forall\,c\in\mathbb{F})\,c\mathbf{u}+\mathbf{v}\in\mathcal{U}$.
\begin{proof}
\leanok
\begin{enumerate}[label=\roman*.]
\setlength{\itemsep}{0pt}
\item[$(\Rightarrow)$] $(\forall\,\mathbf{u},\mathbf{v}\in\mathcal{U})(\forall\,c\in\mathbb{F})\,c\mathbf{u}+\mathbf{v}\in\mathcal{U}$ by inherited scalar and additive closure.
\item[$(\Leftarrow)$] Let $C$ denote the stated closure condition, and choose $\mathbf{u}_0\in\mathcal{U}$. The identity $\mathbf{0}=(-1)\mathbf{u}_0+\mathbf{u}_0\in\mathcal{U}$ holds by $C$ and Corollary \ref{cor:negativeoneactsasinverse}.\\
$(\forall\,\mathbf{u}\in\mathcal{U})\,-\mathbf{u}=(-1)\mathbf{u}+\mathbf{0}\in\mathcal{U}$ by $C$ and Corollary \ref{cor:negativeoneactsasinverse}. $(\forall\,\mathbf{u}\in\mathcal{U})(\forall\,c\in\mathbb{F})\,c\mathbf{u}=c\mathbf{u}+\mathbf{0}\in\mathcal{U}$ by $C$.\\
$(\forall\,\mathbf{u},\mathbf{v}\in\mathcal{U})\,\mathbf{u}+\mathbf{v}=1\mathbf{u}+\mathbf{v}\in\mathcal{U}$ by $C$. The remaining vector-space axioms are inherited from $\mathcal{V}$, so $\mathcal{U}\leq\mathcal{V}$.
\end{enumerate}
\hlt{Structural Commentary}: Nonemptiness supplies a zero; the single closure condition supplies inverses, scalar multiples, and sums.\\
\hlt{Key Insight}: One affine-looking closure test captures all subspace closure.
\end{proof}
\end{theorem}

\begin{proposition} \label{prop:intersectionsubspaces} \hlt{Intersections of Subspaces}\\
Let $\mathcal{V}$ be a vector space over $\mathbb{F}$, let $\emptyset\neq I$, and let $\{\mathcal{U}_i\}_{i\in I}$ be a family of subspaces of $\mathcal{V}$. Then $\bigcap\limits_{i\in I}\mathcal{U}_i\leq\mathcal{V}$.
\begin{proof}
Set $M\defeq\bigcap\limits_{i\in I}\mathcal{U}_i$. Then $\mathbf{0}\in M$ by $(\forall\,i\in I)\,\mathbf{0}\in\mathcal{U}_i$, so $M\neq\emptyset$.\\
$(\forall\,\mathbf{u},\mathbf{v}\in M)(\forall\,c\in\mathbb{F})(\forall\,i\in I)\,c\mathbf{u}+\mathbf{v}\in\mathcal{U}_i$ by the definition of $M$ and $\mathcal{U}_i\leq\mathcal{V}$.\\
Hence $(\forall\,\mathbf{u},\mathbf{v}\in M)(\forall\,c\in\mathbb{F})\,c\mathbf{u}+\mathbf{v}\in M$, and $M\leq\mathcal{V}$ by Theorem \ref{thm:subspacetest}.\\
\hlt{Structural Commentary}: The proof verifies nonemptiness and closure pointwise across the indexed family.\\
\hlt{Key Insight}: Shared closure survives intersection.
\end{proof}
\end{proposition}

\begin{definition} \label{def:linear-combinations-and-span} \hlt{Lists, Linear Combinations, and Span}\\
Let $\mathcal{V}$ be a vector space over $\mathbb{F}$.
\begin{enumerate}[label=(\roman*)]
\setlength{\itemsep}{0pt}
\item For $m\in\N_0$, a \hlt{list} in $\mathcal{V}$ is $\beta=(\mathbf{v}_1,\ldots,\mathbf{v}_m)$; the case $m=0$ is the empty list $()$.
\item For $m\in\N_0$, $\beta=(\mathbf{v}_1,\ldots,\mathbf{v}_m)$, and $(a_1,\ldots,a_m)\in\mathbb{F}^m$, a \hlt{linear combination} of $\beta$ is $\sum\limits_{j=1}^{m}a_j\mathbf{v}_j$.\\
The case $m=0$ is the empty linear combination, whose value is $\mathbf{0}$.
\item For $E\subseteq\mathcal{V}$, its \hlt{span} is the set of all finite linear combinations of entries from $E$.\\
Precisely, $\operatorname{span}(E)\defeq\left\{\sum\limits_{j=1}^{m}a_j\mathbf{v}_j:m\in\N_0,\ (a_1,\ldots,a_m)\in\mathbb{F}^m,\ (\mathbf{v}_1,\ldots,\mathbf{v}_m)\in E^m\right\}$; hence $\operatorname{span}(\emptyset)=\{\mathbf{0}\}$.
\item For $m\in\N_0$ and $\beta=(\mathbf{v}_1,\ldots,\mathbf{v}_m)$, define $\operatorname{span}(\beta)\defeq\operatorname{span}(\{\mathbf{v}_j:1\leq j\leq m\})$. For the empty list, $\operatorname{span}(())=\{\mathbf{0}\}$.
\item $E\subseteq\mathcal{V}$ \hlt{spans} $\mathcal{V}$ iff $\operatorname{span}(E)=\mathcal{V}$. A list $\beta$ \hlt{spans} $\mathcal{V}$ iff $\operatorname{span}(\beta)=\mathcal{V}$.
\end{enumerate}
\end{definition}

\begin{theorem} \label{thm:span-smallest-containing-subspace} \hlt{Span is the Smallest Containing Subspace}\\
Let $\mathcal{V}$ be a vector space over $\mathbb{F}$ and let $E\subseteq\mathcal{V}$. Then:
\begin{enumerate}[label=(\roman*)]
\setlength{\itemsep}{0pt}
\item $\operatorname{span}(E)\leq\mathcal{V}$;
\item $E\subseteq\operatorname{span}(E)$;
\item $(\forall\,\mathcal{U}\leq\mathcal{V})\,[E\subseteq\mathcal{U}\Rightarrow\operatorname{span}(E)\subseteq\mathcal{U}]$.
\end{enumerate}
\begin{proof}
\begin{enumerate}[label=(\roman*)]
\setlength{\itemsep}{0pt}
\item The empty combination has value $\mathbf{0}$, so $\mathbf{0}\in\operatorname{span}(E)$. Let $\mathbf{x},\mathbf{y}\in\operatorname{span}(E)$ and $c\in\mathbb{F}$.\\
Choose $m,n\in\N_0$ and finite linear-combination representations from $E$, $\mathbf{x}=\sum\limits_{j=1}^{m}a_j\mathbf{e}_j$ and $\mathbf{y}=\sum\limits_{k=1}^{n}b_k\mathbf{f}_k$.\\
Then $c\mathbf{x}+\mathbf{y}=\sum\limits_{j=1}^{m}(ca_j)\mathbf{e}_j+\sum\limits_{k=1}^{n}b_k\mathbf{f}_k\in\operatorname{span}(E)$ by concatenating the two lists. Thus $\operatorname{span}(E)\leq\mathcal{V}$ by Theorem \ref{thm:subspacetest}.
\item $(\forall\,\mathbf{v}\in E)\,\mathbf{v}=1\mathbf{v}\in\operatorname{span}(E)$, so $E\subseteq\operatorname{span}(E)$.
\item Let $\mathcal{U}\leq\mathcal{V}$ and $E\subseteq\mathcal{U}$. Set $P(m)\defeq[(\forall\,(a_1,\ldots,a_m)\in\mathbb{F}^m)(\forall\,(\mathbf{e}_1,\ldots,\mathbf{e}_m)\in E^m)\,\sum\limits_{j=1}^{m}a_j\mathbf{e}_j\in\mathcal{U}]$.\\
The base case $P(0)$ holds by $\mathbf{0}\in\mathcal{U}$. Assume $m\in\N_0$, $P(m)$, $(a_1,\ldots,a_{m+1})\in\mathbb{F}^{m+1}$, and $(\mathbf{e}_1,\ldots,\mathbf{e}_{m+1})\in E^{m+1}$.\\
Set $\mathbf{s}\defeq\sum\limits_{j=1}^{m}a_j\mathbf{e}_j$. Then $\mathbf{s}\in\mathcal{U}$ by $P(m)$. $a_{m+1}\mathbf{e}_{m+1}\in\mathcal{U}$ by $E\subseteq\mathcal{U}$ and scalar closure in $\mathcal{U}$.\\
Hence $\sum\limits_{j=1}^{m+1}a_j\mathbf{e}_j=\mathbf{s}+a_{m+1}\mathbf{e}_{m+1}\in\mathcal{U}$ by additive closure, so $P(m+1)$.\\
$(\forall\,m\in\N_0)\,P(m)$ by induction. Hence $\operatorname{span}(E)\subseteq\mathcal{U}$ by the definition of span.
\end{enumerate}
\hlt{Structural Commentary}: The proof verifies closure, generator inclusion, and minimality in statement order.\\
\hlt{Key Insight}: Finite linear combinations are exactly the closure forced by $E$.
\end{proof}
\end{theorem}

\begin{proposition} \label{prop:spanmonotonicity} \hlt{Monotonicity of Span}\\
Let $\mathcal{V}$ be a vector space over $\mathbb{F}$, and let $\mathcal{U},\mathcal{W}\subseteq\mathcal{V}$. Then $\mathcal{U}\subseteq\mathcal{W}\Rightarrow\operatorname{span}(\mathcal{U})\subseteq\operatorname{span}(\mathcal{W})$.
\begin{proof}
Assume $\mathcal{U}\subseteq\mathcal{W}$. Then $\operatorname{span}(\mathcal{W})\leq\mathcal{V}$ and $\mathcal{W}\subseteq\operatorname{span}(\mathcal{W})$ by Theorem \ref{thm:span-smallest-containing-subspace}(i)--(ii).\\
Thus $\mathcal{U}\subseteq\operatorname{span}(\mathcal{W})$, and $\operatorname{span}(\mathcal{U})\subseteq\operatorname{span}(\mathcal{W})$ by Theorem \ref{thm:span-smallest-containing-subspace}(iii).\\
\hlt{Structural Commentary}: The proof compares both generating sets inside the same containing subspace.\\
\hlt{Key Insight}: Minimality converts generator inclusion into span inclusion.
\end{proof}
\end{proposition}

\begin{computationmethod} \label{meth:subspace-diagnostic} \hlt{Subspace and Span Diagnostic}\\
Let $\mathcal{V}$ be a vector space over $\mathbb{F}$ and let $S\subseteq\mathcal{V}$. Return the output matching the task:
\begin{enumerate}[label=(\roman*)]
\setlength{\itemsep}{0pt}
\item \hlt{Subspace certificate}: establish $S\neq\emptyset$ and $(\forall\,\mathbf{u},\mathbf{v}\in S)(\forall\,c\in\mathbb{F})\,c\mathbf{u}+\mathbf{v}\in S$. Then $S\leq\mathcal{V}$ by Theorem \ref{thm:subspacetest}.
\item \hlt{Failure certificate}: if $S=\emptyset$, then $S\nleq\mathcal{V}$. Otherwise exhibit $\mathbf{u},\mathbf{v}\in S$ and $c\in\mathbb{F}$ with $c\mathbf{u}+\mathbf{v}\notin S$.\\
Then $S\nleq\mathcal{V}$ by Theorem \ref{thm:subspacetest}. The condition $\mathbf{0}\in S$ alone is insufficient.
\item \hlt{Generated subspace}: return $\operatorname{span}(S)$, the set of all finite linear combinations from $S$.\\
Precisely, $\operatorname{span}(S)=\left\{\sum\limits_{i=1}^{n}a_i\mathbf{s}_i:n\in\N_0,\ (a_1,\ldots,a_n)\in\mathbb{F}^n,\ (\mathbf{s}_1,\ldots,\mathbf{s}_n)\in S^n\right\}$, with the empty sum $\mathbf{0}$ at $n=0$.
\end{enumerate}
\end{computationmethod}
